% Preliminary manuscript, not the final 8--20-page submission.
\documentclass[11pt]{article}
\usepackage[utf8]{inputenc}
\usepackage[T1]{fontenc}
\usepackage[margin=1in]{geometry}
\usepackage{amsmath,amssymb,amsthm}
\usepackage{booktabs}
\usepackage{xcolor}
\usepackage[round]{natbib}
\usepackage[colorlinks=true,allcolors=blue!50!black]{hyperref}
\newtheorem{proposition}{Proposici\'on}
\newtheorem{lemma}{Lema}
\theoremstyle{definition}
\newtheorem{assumption}{Supuesto}
\title{Reservas bancarias, colateral y asignaci\'on de capital}
\author{Alejandro Ventura\thanks{Universidad del Pac\'ifico. Track B. Repositorio:
\url{https://github.com/alejandroventuraventurameza-tech/ai-project}.}}
\date{Borrador preliminar del n\'ucleo A: 7 de octubre de 2026}
\begin{document}
\maketitle
\begin{abstract}
Una expansi\'on de reservas puede mejorar o empeorar la asignaci\'on de capital.
Este borrador estudia un banco y dos empresas con colateral de pago heterog\'eneo.
Si la empresa productiva est\'a limitada y la otra se autofinancia, reducir el
costo del cr\'edito mejora localmente eficiencia. Si ambas est\'an limitadas,
la reasignaci\'on depende del colateral relativo al patrimonio, y su eficiencia
depende adem\'as de la brecha inicial de productos marginales. Se presentan
derivaciones preliminares y verificaciones num\'ericas. La extensi\'on B, el
ap\'endice manuscrito y la formalizaci\'on Lean siguen pendientes. Este texto
no constituye todav\'ia la entrega final de 8--20 p\'aginas.
\end{abstract}

\section{Introducci\'on}
La pregunta es cu\'ando la liquidez bancaria mejora la asignaci\'on de capital
entre empresas heterog\'eneas. El Track B exige construir el modelo de la tesis;
el alcance A a\'isla un mecanismo est\'atico, con oferta de capital fija. El
alcance B estudiar\'a instrumentos de precio y cantidad, encaje y cambios de
r\'egimen. No se deduce una recomendaci\'on monetaria universal ni bienestar.

\section{Literatura relacionada}
\citet{gonzalez2026}, versi\'on de enero de 2026, es el antecedente principal.
Ya estudia pol\'itica monetaria y mala asignaci\'on. La ec.~4 (p.~7)
impone $q_tk_t\leq\gamma q_ta_t$,
y la ec.~6 (p.~8) expresa el problema est\'atico de la empresa. En esa notaci\'on
$m_t$ es precio real del insumo y $R_t$ costo del capital; son distintos de
nuestras reservas $m$ y precio inicial de alquiler $R$. Sus proposiciones
1--2 (p.~19 y pp.~20--21, pruebas B.7.1--B.7.2) estudian la distribuci\'on
patrimonial y el umbral de productividad. Esta versi\'on tiene tres autores
y reconoce el aporte previo de Albrizio. La relaci\'on general no es nuestra
novedad. Nuestro cierre est\'atico de dos firmas con rendimientos decrecientes
no replica su continuo de firmas con rendimientos constantes, la din\'amica
distributiva ni el problema de pol\'itica \`optima.

\citet{bianchi2022} fundamenta la importancia de reservas, dep\'ositos y
liquidez interbancaria. Se consult\'o localmente su manuscrito de marzo de 2021,
no la versi\'on editorial publicada. Nuestro costo cuadr\'atico es una
aproximaci\'on propia, sin identificarlo con sus ecuaciones. La originalidad
del cierre particular a\'un requiere revisar trabajos citantes.

\section{Modelo}
\begin{assumption}\label{ass:primitives}
$A_H>A_L>0$, $0<\alpha<1$, $K,n_i,h_i>0$, $R>0$ y $t>\rho>0$.
Las primitivas y $K$ no cambian con $t$ en el experimento de asignaci\'on.
\end{assumption}
\begin{assumption}\label{ass:funding}
El capital se alquila al inicio por $Rk_i$, se produce al final
$A_i k_i^\alpha$ y se paga deuda bruta $tb_i$. El patrimonio inicial es $n_i$;
los fondos no utilizados rinden $\rho$. El colateral terminal ex\'ogeno $h_i$
limita el pago, no el principal. Se omite su valor terminal constante.
\end{assumption}
\begin{assumption}\label{ass:clearing}
$k_H+k_L=K$, $B=b_H+b_L$, $B+m=D+E$, $D\geq0$. Un hogar pasivo
ofrece capital y dep\'ositos a los precios dados; no se modela su bienestar.
\end{assumption}
\begin{assumption}\label{ass:bank}
$E\geq0$, $m\geq0$, $c\geq0$, $\kappa>0$, $0<\delta<1$. El banco es
tomador de precios; toma reservas $m$ dadas y elige una cartera interior
factible. Su costo propio es $C=cB+\kappa(\delta D-m)_+^2/2$.
\end{assumption}
La empresa resuelve
\begin{equation}\label{eq:firm}
\max_{k_i,x_i,b_i\geq0} A_i k_i^\alpha-tb_i+\rho(n_i-x_i),
\quad Rk_i=x_i+b_i,\quad x_i\leq n_i,\quad tb_i\leq h_i.
\end{equation}
El banco maximiza $tB+r_m m-\rho D-C$, sujeto a su balance. Sustituir
$D=B+m-E$ da, cuando $\delta D>m$,
\begin{equation}\label{eq:bank}
t=\rho+c+\kappa\delta[\delta(B+m-E)-m].
\end{equation}
La segunda derivada de beneficios en ese interior es $-\kappa\delta^2<0$.
El costo representa escasez de liquidaci\'on de manera reducida; no contiene
b\'usqueda OTC ni una ley de movimiento de reservas. Un encaje activo queda
fuera del experimento A y exige un multiplicador adicional.

\section{Resultados y derivaciones preliminares}
\begin{lemma}\label{lem:funding}
Bajo los supuestos \ref{ass:primitives}--\ref{ass:funding},
$x_i=\min\{Rk_i,n_i\}$ y $b_i=(Rk_i-n_i)_+$. La elecci\'on de capital
es el m\'aximo \`unico de
\[
A_i k_i^\alpha-\rho Rk_i-(t-\rho)(Rk_i-n_i)_++\rho n_i,
\qquad 0\leq k_i\leq(n_i+h_i/t)/R.
\]
\end{lemma}
\begin{proof}
Fijado $k_i$, se tiene $b_i=Rk_i-x_i$. La parte financiera del objetivo es
$-tRk_i+(t-\rho)x_i+\rho n_i$. Como $t>\rho$, aumentar $x_i$ es
estrictamente beneficioso hasta $\min\{Rk_i,n_i\}$. El colateral determina
la cota superior de $k_i$. El t\'ermino productivo es estrictamente c\'oncavo
y el negativo de la parte positiva es c\'oncavo, por lo que la suma es
estrictamente c\'oncava. Continuidad y compacidad garantizan existencia;
concavidad estricta garantiza unicidad. La derivada productiva diverge al
acercarse a cero; con $n_i>0$ la soluci\'on de capital es positiva.
\end{proof}
Escribiendo $M_i=\alpha A_i k_i^{\alpha-1}$, las cuatro regiones son:
$M_i=\rho R$ si $Rk_i<n_i$; $\rho R\leq M_i\leq tR$ en el quiebre;
$M_i=tR$ con deuda interior; y $M_i=tR(1+\mu_i)$, $\mu_i\geq0$,
en la cota de colateral. No se invoca una FOC diferenciable en el quiebre.

\newpage
\begin{proposition}\label{prop:mixed}
Bajo los supuestos \ref{ass:primitives}--\ref{ass:clearing}, sup\'ongase
$b_H=h_H/t$, $M_H>tR$ y $Rk_L<n_L$, de modo que $M_L=\rho R$.
En un intervalo que conserva estas desigualdades, $dk_H/dt<0$.
Una reducci\'on local de $t$ eleva producto y reduce la mala asignaci\'on.
\end{proposition}
\begin{proof}
Si $s=1/(1-\alpha)>1$, el mercado de capital exige
\[
F(R,t)=\frac{n_H+h_H/t}{R}+\left(\frac{\alpha A_L}{\rho R}\right)^s-K=0.
\]
Para cada $t>0$, $F$ es estrictamente decreciente en $R$, diverge a $+\infty$
cuando $R\downarrow0$ y converge a $-K$ cuando $R\uparrow\infty$.
Existe una ra\'iz positiva \`unica. Adem\'as,
\[
F_t=-\frac{h_H}{t^2R},\quad F_R=-\frac{k_H+sk_L}{R}<0,\quad
R'=-\frac{h_H}{t^2(k_H+sk_L)}<0.
\]
Como $k_L=(\alpha A_L/(\rho R))^s$, $k_L'=-sk_LR'/R>0$ y
\[
k_H'=-k_L'=-\frac{sk_Lh_H}{t^2R(k_H+sk_L)}<0.
\]
El producto total cumple $Y'=(M_H-M_L)k_H'<0$, pues $M_H>tR>M_L$.
El benchmark eficiente $k_H^*$ iguala marginales. Con $K$ fijo,
\[
\frac{d}{dk_H}\log\frac{M_H}{M_L}
=(\alpha-1)\left(\frac1{k_H}+\frac1{k_L}\right)<0.
\]
El log-ratio es positivo inicialmente y decrece si aumenta $k_H$; su valor
absoluto tambi\'en decrece mientras no cruce cero. Las desigualdades de
r\'egimen garantizan una mejora local y deben reevaluarse para cambios finitos.
\end{proof}

\newpage
\begin{proposition}\label{prop:both}
Bajo los supuestos \ref{ass:primitives}--\ref{ass:clearing}, sup\'ongase
$b_i=h_i/t$ y $M_i>tR$ para ambas firmas. En ese r\'egimen,
\[
k_H=K\frac{tn_H+h_H}{t(n_H+n_L)+h_H+h_L},\quad
\frac{dk_H}{dt}=\frac{K(n_Hh_L-n_Lh_H)}{[t(n_H+n_L)+h_H+h_L]^2}.
\]
Al caer $t$, aumenta $k_H$ si y solo si $h_H/n_H>h_L/n_L$. Para que ese
aumento mejore eficiencia debe cumplirse $k_H<k_H^*$, y una comparaci\'on
finita no debe sobrepasar $k_H^*$.
\end{proposition}
\begin{proof}
Las dos cotas activas implican $Rk_i=n_i+h_i/t$. Sumando y usando
$k_H+k_L=K$ se obtiene $R=(n_H+n_L+(h_H+h_L)/t)/K$. Sustituir en
la cota de H da la identidad. La regla del cociente produce el numerador
\[
K\{n_H[t(n_H+n_L)+h_H+h_L]-(tn_H+h_H)(n_H+n_L)\}
=K(n_Hh_L-n_Lh_H).
\]
El denominador es positivo; $n_i>0$ permite ordenar los ratios.
Para dos tasas $t_1<t_0$ del mismo r\'egimen, escribiendo $N=n_H+n_L$,
$H_c=h_H+h_L$, se tiene exactamente
\[
k_H(t_1)-k_H(t_0)=
\frac{K(t_1-t_0)(n_Hh_L-n_Lh_H)}{(t_1N+H_c)(t_0N+H_c)}.
\]
La eficiencia requiere adem\'as la brecha de marginales adecuada, pues
$dY/dk_H=M_H-M_L$. La concavidad estricta del producto agregado da
\[
k_H^*=K\frac{A_H^{1/(1-\alpha)}}{A_H^{1/(1-\alpha)}+A_L^{1/(1-\alpha)}}.
\]
Si $k_H<k_H^*$, transferir capital a H sin sobrepasar ese punto aumenta
$Y$ y reduce la brecha absoluta de log-marginales. Si H est\'a por encima,
la misma transferencia empeora ambos indicadores. La igualdad de ratios
de colateral produce participaciones constantes, aunque var\'ie $R$.
\end{proof}

\begin{proposition}\label{prop:reserves}
Bajo el supuesto \ref{ass:bank} y el r\'egimen de la proposici\'on
\ref{prop:mixed} o \ref{prop:both}, con liquidez escasa, balance factible e
interior bancario, aumentar reservas reduce localmente $t$, manteniendo
fijos $\rho,E,c,\kappa,\delta$ y las primitivas reales.
\end{proposition}
\begin{proof}
En el primer r\'egimen $B=h_H/t$ y en el segundo $B=(h_H+h_L)/t$,
por lo que $B'<0$. Diferenciando la ecuaci\'on \eqref{eq:bank} se obtiene
\[
\left(1-\kappa\delta^2B'\right)\frac{dt}{dm}
=-\kappa\delta(1-\delta).
\]
El coeficiente a la izquierda es positivo y el numerador a la derecha es
negativo. La derivada es estrictamente negativa. Con liquidez abundante
el costo marginal adicional es cero y $t=\rho+c$: el efecto local de $m$
es nulo. En el umbral la respuesta debe tratarse por tramos. Si $c=0$,
el tramo abundante permite $t=\rho$ y el lema financiero requiere revisi\'on.
\end{proof}

\section{Verificaci\'on computacional}
El programa \texttt{code/verify.py} comprueba identidades, condiciones KKT
y derivadas mediante soluciones num\'ericas. Usa
10\,000 sorteos, semilla 20261007, y verificaci\'on racional del cambio finito.
En \texttt{code/output/}, los archivos \texttt{verification.json} y
\texttt{allocation\_cases.csv} registran resultados y casos adversos.
El programa falla ante una violaci\'on de las condiciones. Los signos de
eficiencia se comprueban lejos de la igualdad num\'erica de marginales; los
casos neutrales se verifican por separado. Una simulaci\'on no prueba un
teorema universal ni identifica causalidad.

\section{Discusi\'on y limitaciones}
El colateral ex\'ogeno terminal, $K$ fijo y la permanencia de r\'egimen hacen
el trabajo del resultado. Un colateral end\'ogeno al precio del capital,
el encaje activo o la acumulaci\'on patrimonial pueden cambiar los signos.
El costo bancario no sustituye el equilibrio OTC de Bianchi y Bigio. EEA y
ENAHO aportan motivaci\'on descriptiva sin identificar un shock monetario.
Sus proxies y contrafactuales preliminares no son una calibraci\'on auditada.

\section{Conclusi\'on}
La reasignaci\'on causada por una ca\'ida del costo del cr\'edito depende del
r\'egimen financiero y del colateral relativo. Su eficiencia depende adem\'as
de los marginales iniciales. El n\'ucleo A ofrece resultados acotados para
verificaci\'on y una extensi\'on B por desarrollar, sin una recomendaci\'on universal.

\bibliographystyle{plainnat}
\bibliography{references}
\appendix
\section{Pruebas y ap\'endice manuscrito}
Las pruebas preliminares est\'an en el cuerpo. El ap\'endice manuscrito en
\texttt{hand/} est\'a pendiente: la derivaci\'on tipografiada no lo sustituye.

\section{Formalizaci\'on Lean}
Ning\'un resultado tiene todav\'ia una declaraci\'on Lean generada ni un
chequeo completado. Primero se fijar\'a el art\'iculo en un commit integrado.
Se ejecutar\'a AppliedModelingLib sobre ese PDF, con la configuraci\'on del
issue; se conservar\'a la corrida completa y su auditor\'ia.
\begin{center}\small
\begin{tabular}{lll}\toprule
Resultado & Declaraci\'on y archivo & Estado \\\midrule
Lema \ref{lem:funding} & No generados & Abierto: corrida pendiente \\
Proposici\'on \ref{prop:mixed} & No generados & Abierto: corrida pendiente \\
Proposici\'on \ref{prop:both} & No generados & Abierto: corrida pendiente \\
Proposici\'on \ref{prop:reserves} & No generados & Abierto: corrida pendiente \\\bottomrule
\end{tabular}\end{center}
No hay hash formalizado ni salida \texttt{--fast} que reportar. Las
hip\'otesis agregadas y cualquier resultado abierto deben quedar expl\'icitos.

\section{Colaboraci\'on con IA}
Codex redact\'o el cierre y sus derivaciones, implement\'o verificaciones y
prepar\'o LaTeX. Alejandro Ventura a\'un debe revisar los PDF y adjudicar su
contenido. Los prompts originales y las respuestas relevantes est\'an en
\texttt{prompts.md}. La verificaci\'on encontr\'o una exigencia de signo estricto
num\'erico demasiado cerca de igualdad de marginales; se delimit\'o el margen
de precisi\'on y se conservaron casos de neutralidad y signos adversos.
No se registra esta correcci\'on num\'erica como prueba Lean ni veredicto del autor.
\end{document}
